% TOP_PROBLEM: {"id": 30006675, "title": "Global Langlands Functoriality Conjecture", "category_id": 1, "category": {"id": 1, "name": "number_theory", "display_name": "Number Theory", "description": "Properties of integers, prime numbers, Diophantine equations.", "slug": "number-theory", "order_index": 1, "created_at": "2026-07-31T15:26:25.670Z"}, "status": "open"}
% FIRST_POSTED: null
% !TeX program = lualatex
% Compile twice: lualatex open_problems_500.tex
% All references and prose are embedded; no BibTeX database or external inputs.
\documentclass[11pt,a4paper]{article}
\usepackage[margin=24mm,headheight=14pt]{geometry}
\usepackage{fontspec}
\setmainfont{Latin Modern Roman}
\setsansfont{Latin Modern Sans}
\setmonofont{Latin Modern Mono}
\usepackage{amsmath,amssymb,mathtools}
\usepackage{unicode-math}
\setmathfont{Latin Modern Math}
\usepackage{microtype}
\usepackage{enumitem,xurl,needspace}
\usepackage[dvipsnames]{xcolor}
\usepackage{hyperref}
\hypersetup{unicode=true,colorlinks=true,linkcolor=MidnightBlue,urlcolor=MidnightBlue,
 pdftitle={Mathematical Research Notebook},pdfauthor={Source-annotated compilation},bookmarksnumbered=true}
\usepackage{bookmark}
\usepackage{tocloft}
\setlength{\cftsecnumwidth}{3em}
\cftsetpnumwidth{2.5em}
\cftsetrmarg{4em}
\usepackage{titlesec}
\titleformat{\section}{\Large\bfseries\raggedright}{\thesection}{0.7em}{}
\usepackage{fancyhdr}
\pagestyle{fancy}\fancyhf{}
\fancyhead[L]{\small\sffamily Open Problems Math | Research notebook}
\fancyhead[R]{\small 27 September 2026}
\fancyfoot[C]{\thepage}
\setlength{\parindent}{0pt}
\setlength{\parskip}{4pt plus 1pt}
\setlength{\emergencystretch}{3em}
\setcounter{tocdepth}{1}
\setcounter{secnumdepth}{1}
\setlist[itemize]{leftmargin=3em,itemsep=3pt,topsep=4pt,parsep=0pt}
\allowdisplaybreaks
\newcommand{\N}{\mathbb{N}}
\newcommand{\Z}{\mathbb{Z}}
\newcommand{\Q}{\mathbb{Q}}
\newcommand{\R}{\mathbb{R}}
\newcommand{\C}{\mathbb{C}}
\newcommand{\F}{\mathbb{F}}
\newcommand{\A}{\mathbb{A}}
\newcommand{\PP}{\mathbb P}
\newcommand{\E}{\mathbb{E}}
\newcommand{\eps}{\varepsilon}
\newcommand{\dd}{\,\mathrm d}
\newcommand{\sref}[2]{\hyperlink{source:#1:#2}{[S#2]}}
\newcommand{\eref}[2]{\hyperlink{extra:#1:#2}{[E#2]}}
% Turn the next switch off for a shorter reading copy. The quotations preserve
% every exactTarget field, including qualifications omitted from a short summary.
\newif\ifshowcatalogue
\showcataloguetrue
\newcommand{\cataloguescope}[1]{\ifshowcatalogue
 \paragraph{Exact scope in the supplied catalogue.}
 \begin{quote}\small #1\end{quote}\fi}
\usepackage{etoolbox}
\AtBeginEnvironment{itemize}{\small}
% Standalone entry; original section content is delimited for verification.
\begin{document}
\setcounter{section}{0}
%% BEGIN ORIGINAL SECTION CONTAINER
% ================================================================
% problemId: 30006675
% Canonical problem ID: 30006675
% exactTarget SHA-256: a0ece7907c2402fc24072ae6fd7074fecabc7a3cfeaa44b2942779871455e93e
\clearpage
\section*{\texorpdfstring{Global Langlands Functoriality Conjecture}{Global Langlands Functoriality Conjecture}}\label{30006675}
\subsection{Definitions and mathematical statement}
For a number field $F$, $\A_F$ is its ring of adeles. A quasisplit connected reductive group has a Borel subgroup over $F$; its $L$-group incorporates its complex dual group and the Galois action. At an unramified place $v$, the Satake parameter $c(\pi_v)$ is a semisimple conjugacy class in the local $L$-group. For an admissible homomorphism $\rho:{}^LG'\to{}^LG$ over the Galois group, the required weak transfer is
\[
 \forall\pi'\in\operatorname{Aut}(G'),\quad \exists\pi\in\operatorname{Aut}(G),\ \exists S\text{ finite},\quad
 c(\pi_v)=\rho_v\bigl(c(\pi'_v)\bigr)\quad(v\notin S).
\]
Here $\operatorname{Aut}(G)$ denotes automorphic representations in Arthur's inclusive convention, not automorphisms of $G$. The source convention, including admissibility, is retained; no cuspidality of the transfer or uniqueness is imposed.
\par\smallskip\textit{Formulation and concept references:} \sref{30006675}{1}.
\cataloguescope{For every number field F, quasisplit connected reductive F-groups G-prime and G, and admissible L-homomorphism rho:L(G-prime)->L(G) (over Gamma\_F, commuting with its projections, in Arthur Principle1.1's convention), every automorphic representation pi-prime of G-prime(A\_F), in the inclusive Langlands sense cited there, has an automorphic representation pi of G(A\_F) whose unramified Satake classes equal the images under rho\_v of those of pi-prime outside a finite set of places. This is existence of a weak transfer, not uniqueness or a prescribed transfer function.}
\subsection{Short English statement}
Does every allowed map of Langlands $L$-groups produce an automorphic transfer matching the local arithmetic data almost everywhere?
%% BEGIN RESEARCH INSERTION
\subsection{Research attempt}
\paragraph{Current frontier.}
\textit{Research check: 27 September 2026.}
Arthur's Principle 1.1 uses the inclusive automorphic convention, not
only the discrete or cuspidal spectrum. Its reference to Langlands's
1979 paper is substantive: Lemma 1 and Proposition 2 identify global
constituents of parabolic induction from cuspidal data as automorphic
representations. These statements were inspected before using that
convention here. \sref{30006675}{1}, \eref{30006675}{1}.
Established functoriality includes major classical-group transfers and,
for regular algebraic non-CM cuspidal representations of
$\operatorname{GL}_2(\A_{\Q})$, all symmetric powers by Newton--Thorne,
Theorem A. The latter is not a theorem for all number fields and all
admissible maps. Arthur's trace-formula discussion and the 2026 Duke
course outline do not supply such a general theorem. \sref{30006675}{1},
\sref{30006675}{2}, \sref{30006675}{3}, \eref{30006675}{2}.

\paragraph{Attack route.}
A trace-formula comparison needs a global spectral identity isolating
transferred data. A converse-theorem approach needs the appropriate
analytic properties of sufficiently many twisted $L$-functions. A
third route embeds the target dual group in a general linear group,
constructs the composite transfer, and descends it. We first check the
construction in a case where automorphy can actually be supplied, then
stress-test the descent route. None of these approaches may assume an
already constructed universal Langlands group.

\paragraph{Attempt: monomial transfers and an induced subcase.}
Let the source be the split torus $T'=(\mathbf G_m)^r$ over an arbitrary
number field $F$. Its automorphic representations are tuples of
idele-class quasicharacters $\chi_1,\ldots,\chi_r$. First take a split
torus target $T=(\mathbf G_m)^s$ and a dual-torus map given by an integer
matrix $B=(b_{ji})$:
\[
 \rho(z_1,\ldots,z_r)_j=\prod_{i=1}^r z_i^{b_{ji}}.
 \tag{7.1}
\]
Here the $L$-map acts as the identity on the Galois factor and has no
additional Galois cocycle. Define
$\psi_j=\prod_i\chi_i^{b_{ji}}$. Negative exponents cause no problem
because quasicharacters are nonvanishing. Each $\psi_j$ is trivial on
$F^\times$, is continuous, and has the same automorphic smoothness and
moderate-growth properties as a finite product of such characters.
Outside the union of their ramification sets,
\[
 \psi_{j,v}(\varpi_v)=\prod_i\chi_{i,v}(\varpi_v)^{b_{ji}},
 \tag{7.2}
\]
which is precisely the desired equality of torus Satake parameters.
This is an explicit construction in a solved subcase, not a new general
transfer theorem.

Now take target $\operatorname{GL}_s$ and let the restriction of $\rho$
to the complex dual torus be any algebraic $s$-dimensional representation,
still with the trivial cocycle convention. A representation of a complex
torus decomposes into weight spaces, so, after conjugation, its diagonal
weights are (7.1). Put the resulting $\psi_j$ on the diagonal Levi of
the upper-triangular Borel and form normalized parabolic induction. At
an unramified finite place select its spherical irreducible constituent;
at the finitely many remaining places select an irreducible constituent.
The restricted tensor product is automorphic by Langlands's Lemma 1
and Proposition 2. Its unramified Satake class is
\[
 \left[\operatorname{diag}
  \bigl(\psi_{1,v}(\varpi_v),\ldots,\psi_{s,v}(\varpi_v)\bigr)\right].
 \tag{7.3}
\]
Thus it has the required weak transfer. \eref{30006675}{1}.
The normalization is essential: normalized induction uses the square
root of the Borel modulus. Switching induction conventions absorbs this
factor in the inducing torus character; the product formula makes that
modulus character trivial on the rational torus. One must not instead
insert unwanted powers of the residue cardinality into (7.3).
The construction is allowed to be noncuspidal or a subquotient, exactly
as the recorded target permits. Requiring a cuspidal lift would make
this argument inappropriate.

\paragraph{Attempt: composition is valid, but descent is additional.}
Suppose weak transfer is known for admissible maps
${}^LG_0\xrightarrow{\alpha}{}^LG_1\xrightarrow{\beta}{}^LG_2$.
Given $\pi_0$, choose a transfer $\pi_1$, then a transfer $\pi_2$.
Outside the union of their two finite exceptional sets,
\[
 c(\pi_{2,v})=\beta_v(c(\pi_{1,v}))
             =\beta_v\alpha_v(c(\pi_{0,v})).
 \tag{7.4}
\]
This proves transfer for the composite without any uniqueness assumption.
For torus maps it agrees with multiplication of the integer exponent
matrices. An attempted factorization of a general admissible map into
known cases would therefore be legitimate, but no such factorization is
produced here.

Conversely, suppose an admissible map
$r:{}^LG\to{}^L\operatorname{GL}_N$ has been chosen and a transfer for
$r\rho$ constructed. Obtaining a representation $\Pi$ whose Satake
classes equal $r\rho(c(\pi'_v))$ is not yet an automorphic representation
of $G$. The missing assertion must descend \emph{this} global $\Pi$ to
an automorphic representation of $G$ with the prescribed classes
$\rho(c(\pi'_v))$, not merely produce some classes with the same images.
It includes both a global existence issue and a conjugacy issue.

The conjugacy issue already occurs for a faithful embedding of a torus:
\[
 r:\C^\times\hookrightarrow\operatorname{GL}_2(\C),\quad
 r(t)=\begin{pmatrix}t&0\\0&t^{-1}\end{pmatrix},\qquad
 P=\begin{pmatrix}0&1\\1&0\end{pmatrix}.
 \tag{7.5}
\]
We have $Pr(t)P^{-1}=r(t^{-1})$. For $t=2$, the two source elements are
not conjugate in the abelian source, although their images are conjugate
in the target. Thus faithfulness of $r$ does \emph{not} imply injectivity
on semisimple conjugacy classes. Even perfectly matching all characteristic
polynomials after this embedding may lose required local information.
This does not obstruct the explicit torus construction above; it refutes
only a shortcut that tries to reconstruct that construction from the
larger group's conjugacy classes alone.

\paragraph{Stress test.}
The induced subcase uses an established automorphy theorem; merely
forming a restricted tensor product of arbitrary local representations
would not supply it. The general descent argument has no replacement
for that theorem. Strong multiplicity-one reasoning, where available,
can compare representations already known to exist; it cannot create
the missing automorphic representation of $G$. Likewise, equality of
formal Euler products is not an automorphy theorem.

The tensor product (7.3) need not lie in the cuspidal or discrete
$L^2$ spectrum. This is not a defect for the source's inclusive
formulation, but it would be a defect for a silently strengthened
version. The proof has addressed only split torus sources and the
specified untwisted maps; nonsplit groups, nontrivial Galois cocycles
and nonabelian source representations are not covered. No ramified
local prescription, uniqueness, or cuspidality is being claimed.

\paragraph{Outcome.}
\textbf{Outcome: Exploratory approach with unresolved gap.}
The calculations recover an explicit family of valid weak transfers
and prove the composition lemma. The proposed general-linear descent
route has been narrowed to a missing global automorphy statement with
\emph{prescribed} local classes. The faithful-embedding counterexample
shows that an image-only descent assertion is insufficient. Neither
this local correction nor the solved torus subcase resolves the full
functoriality assertion.
%% END RESEARCH INSERTION
\subsection{Sources}
\begin{itemize}\raggedright
\item[\textnormal{[S1]}]\hypertarget{source:30006675:1}{} Classifying automorphic representations, Clay collected-works item76.
\url{https://www.claymath.org/library/cw/arthur/pdf/76.pdf}.
{\small\textit{Catalogue formulation source.}}
\item[\textnormal{[S2]}]\hypertarget{source:30006675:2}{} Functoriality and the Trace Formula — James Arthur.
\url{https://www.math.toronto.edu/arthur/pdf/arthur-func-and-trace-formula.pdf}.
{\small\textit{Catalogue research/status reference; not a proof endorsement.}}
\item[\textnormal{[S3]}]\hypertarget{source:30006675:3}{} Math 590: Open Problems in Number Theory, Spring 2026 — Duke University.
\url{https://sites.math.duke.edu/~pierce/documents/Math590.pdf}.
{\small\textit{Catalogue research/status reference; not a proof endorsement.}}
%% BEGIN RESEARCH REFERENCES
\item[\textnormal{[E1]}]\hypertarget{extra:30006675:1}{} Robert P. Langlands, \emph{On the notion of an automorphic representation}, Proceedings of Symposia in Pure Mathematics 33, Part 1 (1979), 203--207, Lemma 1 and Proposition 2.
\url{https://publications.ias.edu/sites/default/files/On-the-notion-of-an-automorphic-representation-rpl.pdf}.
{\small\textit{Added research source. Primary source incorporated by Arthur's inclusive automorphic convention; induction and irreducible constituents. Consulted 27 September 2026.}}
\item[\textnormal{[E2]}]\hypertarget{extra:30006675:2}{} James Newton and Jack A. Thorne, \emph{Symmetric power functoriality for holomorphic modular forms, II}, arXiv:2009.07180v2 (27 September 2021), Theorem A and Appendix A.
\url{https://arxiv.org/html/2009.07180v2}.
{\small\textit{Added research source. Symmetric powers over Q in the specified holomorphic/regular-algebraic cases, not general functoriality. Consulted 27 September 2026.}}
%% END RESEARCH REFERENCES
\end{itemize}

%% END ORIGINAL SECTION CONTAINER
\end{document}
